\\newpage
\\section{问题重述}

\\subsection{问题背景}
随着"双碳"目标的推进，新型电力系统建设迎来快速发展期。2025年我国可再生能源装机容量已超过煤电，但新能源发电的随机性和波动性给电网调度带来了严峻挑战。根据国家能源局数据，2024年全国弃风弃光率虽降至3.8\\%，但仍存在较大优化空间。

\\subsection{问题提出}
本文针对智能电网负荷预测与多能互补调度优化问题，提出以下三个子问题：

\\textbf{问题一}：基于历史负荷数据和新能源发电数据，分析负荷与发电的时间演化规律，建立高精度的短期负荷预测模型。预测周期为24小时，时间分辨率为1小时。

\\textbf{问题二}：建立多能互补调度环境，定义状态空间、动作空间和奖励函数，设计基于Q-Learning的智能调度Agent，实现电网的实时动态调度。

\\textbf{问题三}：在问题二的基础上，建立多目标优化调度模型，同时考虑调度成本、新能源消纳率和用户满意度三个目标，采用NSGA-II算法求解Pareto前沿，并为决策者提供最优调度方案。

\\subsection{研究意义}
本研究对于提升新型电力系统的运行效率、促进新能源消纳、保障电网安全稳定运行具有重要的理论价值和工程应用前景。
